\documentclass[11pt,a4paper]{article}
\usepackage[T1]{fontenc}
\usepackage{lmodern,amsmath,amssymb,amsthm,mathtools}
\usepackage[margin=25mm]{geometry}
\usepackage{microtype,booktabs}
\usepackage[hidelinks]{hyperref}
\hypersetup{pdftitle={Linear Asymptotics and a Disk-Smearing Bound for the Mordell Distance Product},pdfauthor={Fengcan Jiang}}
\setlength{\emergencystretch}{2em}
\newtheorem{theorem}{Theorem}[section]
\newtheorem{proposition}[theorem]{Proposition}
\newtheorem{lemma}[theorem]{Lemma}
\newtheorem{corollary}[theorem]{Corollary}
\theoremstyle{remark}
\newtheorem{remark}[theorem]{Remark}
\newcommand{\C}{\mathbb C}
\newcommand{\R}{\mathbb R}
\newcommand{\dd}{\,d}
\newcommand{\ind}{\mathbf 1}
\newcommand{\D}{\mathcal D}
\newcommand{\pos}[1]{\left(#1\right)_+}
\numberwithin{equation}{section}
\title{Linear Asymptotics and a Disk-Smearing Bound\\
for the Mordell Distance Product}
\author{Fengcan Jiang}
\date{27 October 2026}
\begin{document}
\maketitle
\begin{abstract}
For the maximum \(M_n\) of the squared distance product under
\(\sum_{i=1}^n|z_i|^2=n\), we give an exact finite-\(n\) conversion to the
ground-state energy of a planar Coulomb gas with quadratic confinement.
The deterministic theorem of Sandier and Serfaty then yields
\(\log M_n=(\frac12\log2-\frac14)n^2+\frac n2\log n+\gamma n+o(n)\).
The microscopic constant satisfies
\(0.4021788952233<\gamma<0.4034264097201\).
The lower endpoint is the triangular-lattice comparison of
B\'etermin and Sandier; the upper endpoint agrees with the classical
disk-smearing bound discussed by Lauritsen.
A self-contained smearing argument also proves an explicit finite-\(n\)
upper bound. Its exact deficit splits into Coulomb discrepancy,
boundary leakage, and pair-overlap terms. For eight points this gives a
concrete positive-gap problem, which remains unsolved.
\end{abstract}

\section{The variational problem and its normalization}
All logarithms are natural. For \(n\ge2\), put \(N=n(n-1)/2\) and define
\begin{equation}\label{eq:Mn}
 \Delta_n(z)=\prod_{i<j}|z_i-z_j|^2,\qquad
 M_n=\max_{\substack{z\in\C^n\\\sum_i|z_i|^2=n}}\Delta_n(z).
\end{equation}
Compactness gives a maximizer, and a configuration of distinct points
shows that its value is positive. In particular, a maximizer has no
collisions. We use the ordered-pair Coulomb Hamiltonian
\begin{equation}\label{eq:H}
 H_n(x)=-\sum_{i\ne j}\log|x_i-x_j|
              +\frac n2\sum_i|x_i|^2,\qquad E_n=\min_{\C^n}H_n .
\end{equation}
Thus \(H_n=-\log\Delta_n+(n/2)Q\), where \(Q=\sum_i|x_i|^2\).
The confinement controls the logarithmic interaction at infinity, and
collisions have infinite energy, so the minimum is attained away from
collisions.

\begin{proposition}[Exact global conversion]\label{prop:bridge}
For every \(n\ge2\),
\begin{equation}\label{eq:bridge}
 \log M_n=A_n-E_n,\qquad
 A_n=N\left[1-\log\left(1-\frac1n\right)\right].
\end{equation}
Every minimizer of \(H_n\) has \(Q=n-1\). Multiplication by
\(\sqrt{n/(n-1)}\) maps such minimizers to maximizers in \eqref{eq:Mn},
and the inverse scaling maps maximizers to minimizers.
\end{proposition}
\begin{proof}
For every \(Q>0\), homogeneity gives
\[
 \max_{\sum_i|x_i|^2=Q}\Delta_n(x)=(Q/n)^N M_n.
\]
Consequently \(E_n\) is the minimum, over \(Q>0\), of
\[
 -\log M_n-N\log(Q/n)+nQ/2.
\]
Its derivative is \(-N/Q+n/2\), and its second derivative is \(N/Q^2>0\).
Its unique minimum is \(Q=2N/n=n-1\). Evaluation proves
\eqref{eq:bridge} and the assertions about equality and scaling.
\end{proof}
This is a global identity; no prescribed shape or stationary-point
classification enters it. Its normalization correction has the
convergent expansion
\begin{equation}\label{eq:Aseries}
 A_n-\frac{n^2}{2}
 =-\frac12\sum_{k=0}^{\infty}
       \frac{n^{-k}}{(k+1)(k+2)}
 =-\frac14-\frac1{12n}+O(n^{-2}).
\end{equation}
Keeping the first tail term and bounding the rest geometrically yields
\[
 -\frac14-\frac1{12(n-1)}
 \ \le\ A_n-\frac{n^2}{2}\
 \le-\frac14-\frac1{12n}.
\]
In particular, the conversion contributes no additional term of order
\(n\). The constant \(-1/4\) in this conversion is not an established
constant term in the asymptotic expansion of \(\log M_n\).

\section{The Coulomb theorem and the coefficient gamma}
For \(g(x)=-\log|x|\), write
\[
 \D(\mu)=\iint g(x-y)\dd\mu(x)\dd\mu(y),\qquad
 I_V(\mu)=\D(\mu)+\int V\dd\mu,\quad V(x)=|x|^2/2.
\]
There is no factor \(1/2\) in \(\D\). The equilibrium probability measure
and its value are
\begin{equation}\label{eq:equilibrium}
 \mu_0=\frac1{2\pi}\ind_{\{|x|\le\sqrt2\}}\dd x,\qquad
 I_0=I_V(\mu_0)=\frac34-\frac12\log2.
\end{equation}
Indeed, the logarithmic potential of a uniform disk of radius \(R\)
equals \(\log R-\frac12+|x|^2/(2R^2)\) inside the disk and
\(\log|x|\) outside it. For \(R=\sqrt2\), its negative plus \(V/2\) is
constant on the disk and larger outside. The disk's second moment is
one, giving \eqref{eq:equilibrium}.

We specify the energy convention needed for the linear coefficient.
Let \(\mathcal A_1\) be the class in Sandier--Serfaty
\cite{SS}: its fields satisfy
\[
 \operatorname{div}E=2\pi(\nu-1),\qquad
 \operatorname{curl}E=0,\qquad
 \nu=\sum_{p\in\Lambda}\delta_p,
\]
where \(\Lambda\) is discrete and
\(\sup_{R>1}\nu(B_R)/|B_R|<\infty\).
With the cutoffs \(\chi_R\) of their Definition 1.3, set
\[
 W(E)=\limsup_{R\to\infty}\frac1{|[-R,R]^2|}
 \lim_{\eta\downarrow0}\left[
 \frac12\int_{\R^2\setminus\bigcup_{p\in\Lambda}B(p,\eta)}
                  \chi_R|E|^2
       +\pi\log\eta\sum_{p\in\Lambda}\chi_R(p)\right].
\]
We write \(W_*=\min_{\mathcal A_1}W\).
The choice of density and of the half-Dirichlet normalization is part
of this definition. A different convention changes the numerical
value assigned to \(W_*\).

\begin{theorem}[Linear asymptotic term]\label{thm:asymptotic}
There is a constant, independent of \(n\),
\begin{equation}\label{eq:gamma}
 \gamma=-\frac{W_*}{\pi}-\frac12\log(2\pi)
\end{equation}
such that
\begin{equation}\label{eq:asymptotic}
 \log M_n=\left(\frac12\log2-\frac14\right)n^2
                 +\frac n2\log n+\gamma n+o(n).
\end{equation}
\end{theorem}
\begin{proof}
The quadratic potential satisfies the growth and regularity
hypotheses (1.11)--(1.13) of \cite{SS}: it is smooth,
\(V/2-\log|x|\to+\infty\), its Laplacian is the positive constant two,
and the equilibrium support has smooth circular boundary.
The deterministic conclusion of Theorem 2(C), together with their
equations (1.23) and (1.40), therefore gives
\[
 E_n=n^2I_0-\frac n2\log n+
 \left[\frac{W_*}{\pi}-\frac12\int m_0\log m_0\dd x\right]n+o(n).
\]
Here \(m_0=1/(2\pi)\) on its support and has total mass one, so the
coefficient in brackets is \(W_*/\pi+\frac12\log(2\pi)\).
Substitution into \eqref{eq:bridge}, using \eqref{eq:Aseries} and
\eqref{eq:equilibrium}, proves the claim.
\end{proof}
The proof applies the deterministic ground-state theorem directly.
No exchange of the large-\(n\) and zero-temperature limits is required.
This use of the published theorem, also formulated in
B\'etermin--Sandier \cite{BS}, is an external analytic input.

\section{A finite-size disk-smearing bound}
Let \(\sigma_r\) be uniform probability measure on the disk \(B_r\),
and let \(T_rf=f*\sigma_r\). Direct angular averaging gives
\begin{equation}\label{eq:potential}
 (T_rg)(d)=
 \begin{cases}
 -\log r+\frac12-\dfrac{d^2}{2r^2},&0\le d<r,\\
 -\log d,&d\ge r.
 \end{cases}
\end{equation}
In particular \(T_rg\le g\). Positivity of the averaging operator then
gives \(T_r^2g\le T_rg\le g\). These inequalities allow overlapping
disks; no point-separation assumption is made.
Integration of \eqref{eq:potential} gives the two constants
\begin{equation}\label{eq:self}
 \D(\sigma_r)=-\log r+\frac14,\qquad
 \int|u|^2\dd\sigma_r(u)=\frac{r^2}{2}.
\end{equation}

\begin{theorem}[Explicit global upper bound]\label{thm:finite}
For every \(n\ge2\),
\begin{equation}\label{eq:Un}
 M_n\le U_n:=
 n^{n/2}\left(\frac{2n}{n-1}\right)^N e^{-N/2}.
\end{equation}
Consequently
\begin{equation}\label{eq:gammaU}
 \gamma\le\gamma_D:=\frac34-\frac12\log2.
\end{equation}
\end{theorem}
\begin{proof}
For any collision-free points \(x_i\), set
\(\mu_x=n^{-1}\sum_i\sigma_r(\,\cdot-x_i)\).
Including and then subtracting exactly \(n\) self-terms yields
\[
 \sum_{i\ne j}g(x_i-x_j)
 \ge n^2\D(\mu_x)-n\D(\sigma_r).
\]
Moreover,
\(n^2\int V\dd\mu_x=(n/2)\sum_i|x_i|^2+n^2r^2/4\).
Since \(I_V(\mu_x)\ge I_0\), \eqref{eq:self} gives
\[
 H_n(x)\ge n^2I_0+n\log r-\frac n4-\frac{n^2r^2}{4}.
\]
The right side is maximal at \(r^2=2/n\). Thus
\begin{equation}\label{eq:Ln}
 E_n\ge L_n:=
 n^2I_0-\frac n2\log n
           +n\left(\frac12\log2-\frac34\right).
\end{equation}
The exact conversion \eqref{eq:bridge} simplifies to
\(A_n-L_n=\log U_n\), proving \eqref{eq:Un}.
Comparison with \eqref{eq:asymptotic} proves \eqref{eq:gammaU}.
\end{proof}
This finite-\(n\) proof does not use the asymptotic theorem.
For example, \eqref{eq:Aseries} also gives the strict bound
\[
 \log M_n<
 \left(\frac12\log2-\frac14\right)n^2+\frac n2\log n
                    +\gamma_Dn-\frac14.
\]
The final \(-1/4\) is a correction in this upper bound, not an established
boundary correction for the true optimum.

The upper endpoint \eqref{eq:gammaU} agrees with the classical
Lieb--Narnhofer and Sari--Merlini smearing bound discussed in
Lauritsen \cite{Lau}; see Appendix B and Remark 4.5 in the cited
arXiv version. We do not claim this universal endpoint as a new
constant bound. The preceding argument gives its direct specialization
to the finite moment-constrained distance product.

\section{The triangular comparison and numerical certificates}
B\'etermin--Sandier \cite[Section 7.3]{BS} compute the energy of the
unit-density triangular lattice:
\[
 \frac1\pi W(\Lambda_\triangle)
  =\log\pi-\frac12\log3-3\log\Gamma(1/3)+\frac32\log2 .
\]
Because \(W_*\le W(\Lambda_\triangle)\), the negative sign in
\eqref{eq:gamma} gives
\begin{equation}\label{eq:interval}
 \gamma_\triangle\le\gamma\le\gamma_D,\qquad
 \gamma_\triangle=
 3\log\Gamma(1/3)+\frac12\log3-2\log2-\frac32\log\pi.
\end{equation}
The companion rational interval certificates give, in particular,
\[
 \boxed{\quad
 0.4021788952233<\gamma<0.4034264097201 .
 \quad}
\]
These decimal endpoints strictly enclose the constants in
\eqref{eq:interval}; they are not a numerical identification of
\(\gamma\). Identifying \(\gamma\) with \(\gamma_\triangle\) would require
the triangular value to minimize \(W\) over the full admissible class.
Optimality within the Bravais-lattice family alone is insufficient.

The program \texttt{verify\_gamma\_bounds.py} certifies the numerical
enclosures using rational arithmetic, a Machin formula for \(\pi\),
integer square roots, and controlled exponential, logarithmic and
Dedekind-eta series. It evaluates an established analytic lattice
formula; it does not prove the Coulomb asymptotic theorem or the
global crystallization assertion.

\section{An exact nonnegative deficit and the eight-point problem}
Keep \(r^2=2/n\), let \(K_r=T_r^2g\), and define
\begin{equation}\label{eq:zeta}
 \zeta(\rho)=
 \begin{cases}
 0,&0\le\rho\le\sqrt2,\\
 \rho^2/4-\log\rho-\frac12+\frac12\log2,&\rho>\sqrt2.
 \end{cases}
\end{equation}
This is the equilibrium obstacle excess. It is nonnegative.
For \(\nu_x=\mu_x-\mu_0\), expansion about \(\mu_0\) gives
\begin{equation}\label{eq:meanfielddeficit}
 I_V(\mu_x)-I_0=\D(\nu_x)+2\int\zeta\dd\mu_x.
\end{equation}
The signed measure \(\nu_x\) has zero total mass, bounded density and
compact support. If \(U_{\nu_x}=g*\nu_x\), integration by parts gives
\[
 \D(\nu_x)=\frac1{2\pi}\int_{\R^2}|\nabla U_{\nu_x}|^2\dd x\ge0.
\]
The zero mass eliminates the boundary term at infinity.

\begin{proposition}[Exact deficit identity]\label{prop:deficit}
Every collision-free configuration satisfies
\begin{equation}\label{eq:deficit}
 \begin{split}
 H_n(x)-L_n={}&n^2\D(\nu_x)+2n^2\int\zeta\dd\mu_x\\
 &+\sum_{i\ne j}
      \bigl[g(x_i-x_j)-K_r(x_i-x_j)\bigr].
 \end{split}
\end{equation}
All three terms on the right are nonnegative. The point-pair term
vanishes for a pair at distance at least \(2r\).
\end{proposition}
\begin{proof}
The exact smoothed-pair sum, including self-terms, is
\(n^2\D(\mu_x)=\sum_{i\ne j}K_r(x_i-x_j)+n\D(\sigma_r)\).
Combine this with the exact smoothed moment formula in the proof of
Theorem~\ref{thm:finite}, \eqref{eq:self}, and \(r^2=2/n\).
The result is
\[
 H_n-L_n=n^2[I_V(\mu_x)-I_0]+\sum_{i\ne j}(g-K_r)(x_i-x_j).
\]
Now use \eqref{eq:meanfielddeficit}. The kernel comparison proves
nonnegativity. If \(d\ge2r\), the first smoothed potential agrees with
\(g\) throughout the second disk, and its second average also equals
\(g(d)\); hence the pair gap is zero.
\end{proof}
There is also a simple explicit lower estimate. For \(d>0\),
\begin{equation}\label{eq:pairlower}
 g(d)-K_r(d)\ge
 \pos{1-d/r}^{\,2}+\frac1{16}\pos{1-d/(2r)}^{\,4}.
\end{equation}
To see this, put \(h_r=g-T_rg\). Formula~\eqref{eq:potential} gives
\(h_r(d)\ge\pos{1-d/r}^2\), and \(g-K_r=h_r+T_rh_r\).
If \(d<2r\), put \(\varepsilon=1-d/(2r)\).
The disk of radius \(r\varepsilon/2\) at the midpoint of the two
centers lies in their overlap. Its normalized area is
\(\varepsilon^2/4\), and \(h_r\ge\varepsilon^2/4\) there.
This proves the second term of \eqref{eq:pairlower}.

For \(Q(x)=n-1\), define \(z=\sqrt{n/(n-1)}x\). The conversion
\eqref{eq:bridge} now holds for this individual configuration as
\begin{equation}\label{eq:configdeficit}
 H_n(x)-L_n=\log U_n-\log\Delta_n(z).
\end{equation}
For eight points the conjectured sharp value is
\(C_8=8^{28}/7^{21}\). Thus its global upper bound is equivalent to
the following positive-gap assertion over all configurations with
\(Q(x)=7\):
\begin{equation}\label{eq:gap8}
 H_8(x)-L_8\ge\delta_8,\qquad
 \delta_8=\log(U_8/C_8)=40\log2-7\log7-14 .
\end{equation}
Here the smoothing radius is \(r=1/2\). The rational companion
certificate establishes, for example,
\[
 0<\delta_8<0.104517,\qquad 1<U_8/C_8<1.110174.
\]
Equation~\eqref{eq:deficit} therefore reformulates the remaining
eight-point problem in terms of three continuous, nonnegative costs.
It does not establish the positive gap in \eqref{eq:gap8}, and does not
exclude the general nonzero-moment branch by itself.

\section{Verification, attribution and scope}
The companion archive contains two standard-library acceptance programs:
\path{verify_gamma_bounds.py} and \path{verify_finite_bound.py}.
They include rational interval certificates. A unified entry point
checks source hashes and runs them in a temporary copy:
\begin{center}
 \texttt{python -B -S verify\_all.py}.
\end{center}
The programs audit constants, algebraic bookkeeping and interval
enclosures. The analytic averaging inequalities and the external
ground-state theorem remain parts of the mathematical proof; this is
not a proof-assistant formalization.

The linear coefficient exists by the published theorem of
Sandier--Serfaty, its triangular comparison uses
B\'etermin--Sandier, and the strong upper endpoint is consistent with
the classical smearing literature summarized by Lauritsen.
This paper records a normalization audit, a finite-\(n\) specialization,
and an exact deficit formulation. It makes no global priority claim.

The project and source material are attributed to Fengcan Jiang.
Preparation, computation and cross-checking included assistance from
OpenAI Codex. AI-assisted checking is not independent external peer
review. The exact microscopic minimum, an effective uniform rate in the
\(o(n)\) remainder, true finite-\(n\) boundary corrections, and the
unconditional eight-point sharp inequality are not established here.

\begin{thebibliography}{9}
\bibitem{SS}
E.~Sandier and S.~Serfaty,
\emph{2D Coulomb Gases and the Renormalized Energy},
arXiv:1201.3503v2 (2013).
\url{https://arxiv.org/abs/1201.3503v2}.

\bibitem{BS}
L.~B\'etermin and E.~Sandier,
\emph{Renormalized Energy and Asymptotic Expansion of Optimal
Logarithmic Energy on the Sphere},
Constructive Approximation \textbf{47} (2018), 39--74;
arXiv:1404.4485v5.
\url{https://doi.org/10.1007/s00365-016-9357-z}.

\bibitem{Lau}
A.~B.~Lauritsen,
\emph{Floating Wigner crystal and periodic jellium configurations},
Journal of Mathematical Physics \textbf{62} (2021), 083305;
arXiv:2103.07975v2.
\url{https://doi.org/10.1063/5.0053494}.
\end{thebibliography}
\end{document}
